Eine \mSbO\ schwere Snowboarderin fährt mit einer Geschwindigkeit von \vSbO\ durch eine Senke und anschliessend über einen Buckel mit je \rSbO\ Krümmungsradius.

\begin{abcliste}
\abc Zeichnen Sie die Kräfte auf die Boarderin im tiefsten Punkt der Senke ein. Aus welchen Kräften setzt sich die Zentripetalkraft zusammen?
\abc Wie gross ist die Normalkraft vom Boden auf das Board im tiefsten Punkt der Senke bzw. im höchsten Punkt des Buckels?

Hinweis: Stellen Sie einen Ausdruck für die resultierende Kraft auf, und benutzen Sie die Zentripetalbedingung.
\abc Wie schnell kann die Snowboarderin über den Buckel fahren, ohne dass sie abhebt? Wie muss sie sich verhalten, damit dies auch mit höheren Geschwindigkeiten möglich ist?
\end{abcliste}
